\section{Why this matters for cryptographic and decentralised systems}
\label{sec:crypto}

The second constituency wants the same four properties as the first, for
different reasons. An agent wants canonical, addressable, verifiable,
attributable knowledge because it cannot trust its own memory. A decentralised
system wants it because there is no server to trust.

\subsection{The gap content addressing leaves}

Content-addressed storage --- Merkle structures \citep{merkle1987}, IPFS
\citep{benet2014}, Ethereum Swarm \citep{tron2020} --- solved integrity and
made durable, host-independent references possible. A hash names exactly one
immutable byte string, forever, and tampering is detectable on retrieval.

What it did not solve is \emph{meaning}. A 32-byte reference tells you nothing
about what the content is or how it relates to anything else. The standard
answer has been the manifest: a mapping from path names to references --- in
other words, a directory tree. And a directory tree is the wrong shape for the
same reason it was always the wrong shape: every object must live in exactly
one place, while almost everything worth storing is legitimately several
things at once. \texttt{outbound.pdf} is a flight document, a Japan document,
and a refundable booking; a tree makes you choose, or makes you duplicate.

\begin{figure}[h]
\centering
\begin{subfigure}[b]{0.47\textwidth}\centering
  \dotfig[0.82]{folder_tree}
  \caption{A tree must choose. The Paris flight is under
  \textsf{Flights}; the Japan one is under \textsf{Japan} --- so
  ``all flights'' needs a duplicate.}
\end{subfigure}\hfill
\begin{subfigure}[b]{0.47\textwidth}\centering
  \dotfig[0.82]{dag_multiparent}
  \caption{A DAG does not. \texttt{outbound.pdf} is under both, once, and
  the query decides which view you get.}
\end{subfigure}
\caption{The manifest problem in one picture.}
\label{fig:tree-vs-dag}
\end{figure}

\odag{} is a semantic layer for exactly this gap: a small structure that
records \emph{what things are}, independently of where their bytes live, and
that happens to have the one property content-addressed storage rewards above
all others.

\subsection{Semantic content addressing}

\begin{keyidea}
Ordinary content addressing gives \textbf{same bytes, same hash}. \odag{}
gives \textbf{same meaning, same hash}.
\end{keyidea}

That upgrade is the whole contribution, and it is worth spelling out what it
takes. Two people file the same six documents. They do it in different orders.
One of them types a few redundant links along the way. One writes
$400\,\mathrm{g}$ and the other $0.4\,\mathrm{kg}$. Ordinary content
addressing gives them two different hashes, because their bytes differ.
\odag{} gives them one, because:

\begin{itemize}[leftmargin=1.4em]
\item the transitive reduction is unique, so build order and redundant edges
  are quotiented away (Thm~\ref{thm:unique-reduction});
\item values are canonicalised by denotation, so spelling is quotiented away;
\item the serialisation and the storage trie are canonically encoded, so
  nothing about \emph{how} the store was written survives into the bytes.
\end{itemize}

The consequences are unusual enough to list, because no mainstream database
offers them:

\begin{enumerate}[leftmargin=1.7em]
\item \textbf{Deduplication of meaning.} Two organisations that independently
  build the same taxonomy converge on one address, and discover it by
  comparing 32 bytes.
\item \textbf{Agreement is a hash comparison.} ``Do we share an ontology?''
  costs one string comparison and zero data transfer. Because canonicality is
  semantic, this is same-\emph{meaning}-same-hash, not
  same-\emph{bytes}-same-hash.
\item \textbf{A root is a commitment.} Publishing a root commits you to an
  entire knowledge state, in the cryptographic sense: you cannot later change
  what you knew without changing the reference. This is a commitment scheme
  you get for free, without designing one.
\item \textbf{Snapshots are free and eternal.} Nothing is overwritten, so
  every past root still resolves. ``What did the catalogue say on the day of
  the contract?'' is answerable years later.
\item \textbf{Disagreement localises.} Two differing roots can be narrowed to
  the exact differing records by structural diff, without transferring either
  store.
\end{enumerate}

\subsection{Convergence without consensus}

This is the property that makes \odag{} suit decentralised infrastructure
rather than merely tolerate it, and it deserves emphasis because it is
frequently assumed to be impossible.

Distributed systems usually achieve agreement by \emph{consensus}: a protocol
(or a blockchain) that decides an order of operations. Consensus is expensive,
requires liveness, and requires participants to be online together.

\odag{} needs none of it. Merge is commutative and idempotent
(Thm~\ref{thm:merge}), so replicas that have seen the same set of updates ---
in any order, with any duplicates --- hold identical state. This is the
defining property of a state-based CRDT \citep{shapiro2011}, and combined with
content addressing it gives the pattern \citet{sanjuan2020} call a
Merkle-CRDT: state identified by hash, merged by a convergent function,
gossiped rather than ordered, and which is the substrate the local-first
software argument \citep{kleppmann2019} asks for: full local capability,
collaboration without a coordinating server, and no vendor holding the
authoritative copy.

\begin{keyidea}
Two writers who fold in each other's published roots reach the
\emph{byte-identical} root regardless of gossip order (G5). No server, no
chain, no lock, no coordination --- and afterwards they can \emph{prove} they
agree by comparing the root.

The blockchain, when one is involved at all, is therefore not needed for
agreement. It is needed only for the things chains are actually good at:
timestamping, settlement, and adjudicating disputes about money.
\end{keyidea}

The price is stated in Section~\ref{sec:merge} and repeated here because it is
the one thing an integrator must design around: merge is union, so it is
grow-only. A removal can be resurrected by a peer holding the old state.
Undoing is local. Any application that needs authoritative deletion must add
its own layer.

\subsection{Publishing, following, and paying rent}
\label{sec:crypto-economics}

Concretely, on Swarm \citep{tron2020}, an \odag{} store becomes:

\begin{center}\small
\begin{tabularx}{\textwidth}{@{}p{4.2cm} X@{}}
\toprule
\textbf{Swarm mechanism} & \textbf{What \odag{} uses it for}\\
\midrule
chunks, retrieved by hash & the record blobs and trie nodes: one commit
  writes only what changed\\
postage stamps (storage rent) & keeping a published ontology alive; and, more
  interestingly, \emph{not} keeping derived indexes alive\\
signed feeds & a followable address for ``the latest root published by this
  key'' --- mutability on top of immutability\\
retrievability checks & confirming a published ontology is actually fetchable
  by strangers\\
\bottomrule
\end{tabularx}
\end{center}

The postage mechanism produces a genuinely elegant fit with the derived-layer
discipline of Section~\ref{sec:keep}. Derived structures --- cone summaries,
implication lints, MDL-learned concepts, materialised meets --- are
regenerable by definition. On a network where storage costs rent, that gives
eviction a native mechanism:

\begin{keyidea}
\textbf{Let a derived index's postage lapse.} A cache whose eviction policy is
``stop paying'' is a cache with an honest cost model: an index is kept warm
exactly as long as query traffic justifies its rent, and its disappearance
costs nothing but recomputation. Regenerable things are allowed to die.
\end{keyidea}

This also gives the workload-optimal materialisation problem of
Section~\ref{sec:dir-materialisation} a real objective function rather than a
notional one: the space term is not a tuning constant, it is money.

\subsection{Anchoring, and settling disputes}

A root is 32 bytes, which is exactly the size of one storage word in the
Ethereum Virtual Machine. Putting a root in a contract or an event log is a
timestamped commitment to a whole knowledge state, at negligible cost.

What that enables is dispute resolution by proof rather than by argument:

\begin{enumerate}[leftmargin=1.7em]
\item A party publishes a catalogue and anchors its root on chain.
\item A counterparty acts on a claim about the catalogue --- ``this item is
  under $5\,\mathrm{kg}$ and located in the EU''.
\item If the claim is disputed, it is settled by exhibiting a proof against
  the anchored root. Inclusion, absence and subsumption all have proofs
  (\S\ref{sec:agents-verify}).
\end{enumerate}

Two practical tailwinds are worth recording. Swarm's addressing is
keccak-based --- the EVM's native hash --- so on-chain verification of a
Swarm-side inclusion proof does not require an expensive foreign hash
function; and Swarm's own storage-incentive machinery already verifies such
proofs on chain, so the pieces exist.

One honest gap, though. \odag{}'s current subsumption certificates work by
\emph{re-execution over authenticated records}, which is ideal off-chain
(semantics stay single-sourced in the real implementation) and impractical
on-chain (a contract cannot re-run a graph traversal). The shape that a
contract \emph{can} check is a hash-link path: a Merkle-ised index over the
semantic graph, where a positive subsumption is a self-verifying chain of
hashes. That is a known, designed-but-unbuilt extension, and it is the right
thing to build when --- and only when --- somebody actually needs a
subsumption verified on chain.

\subsection{Bonded claims: the economic third leg}
\label{sec:crypto-bonds}

Limit L1 says proofs attest structure, never truth. Attribution
(Section~\ref{sec:agents}) upgrades that to ``someone said it, and you may
choose whose word you take''. There is a third leg available, and it is
economic: \emph{someone will pay if it is wrong}.

A sister project (\texttt{factbond}, design stage) explores bonded assertions:
a stake slashed on successful dispute, a premium that prices reliability.
Whatever the details, two structural gifts flow from \odag{}'s side of the
fence and are worth naming because they are unusual:

\begin{itemize}[leftmargin=1.4em]
\item \textbf{Claim identity is free and canonical.} Bonding a claim requires
  naming it unambiguously. A root plus a canonical claim is exactly that ---
  and because a root commits to the whole store, a single bond can cover
  millions of facts, with any individual one disputed later by an inclusion
  proof. \emph{Batched commitment, itemised dispute.}
\item \textbf{One class of dispute adjudicates mechanically.} Distinguish
  ``matches-source'' claims (``the store at root $R$ entails $X$'') from
  ``matches-world'' claims (``$X$ is true''). The first is settled by a Tier-2
  certificate --- a proof checker is the cheapest possible rung of a dispute
  ladder, and it needs no human and no oracle. Only matches-world claims need
  evidence and adjudicators.
\end{itemize}

That split is only available because subsumption is decidable from the stored
graph plus exact arithmetic. It is a direct dividend of the refusals in
Section~\ref{sec:keep}.

\subsection{Marketplaces: matching as cone intersection}

A second sister project (\texttt{loopmarket}) uses \odag{} for marketplace
matching, and it is the application that forced typed values into existence.
A courier's constraint --- ``anything up to $5\,\mathrm{kg}$, collected in
this region, delivered before Friday'' --- must match offers nobody
pre-classified against that threshold.

Three properties of the result matter for a crypto audience:

\begin{itemize}[leftmargin=1.4em]
\item \textbf{Matching is the query primitive.} A constraint is a query term,
  so matching is one cone intersection --- one planner, one merge story, no
  second filtering mechanism beside subsumption.
\item \textbf{Recall-completeness is guaranteed, not hoped for.}
  \cmd{get\_overlapping} returns every value whose denotation provably
  intersects the constraint's (G6). For a marketplace, ``we did not miss a
  match'' is a property worth being able to state.
\item \textbf{No floats, anywhere.} Values are exact rationals of an SI anchor
  unit. This is the same discipline crypto already uses for money --- wei,
  satoshi, integer base units --- adopted for the same reason: exactness by
  construction rather than by careful rounding.
\end{itemize}

\subsection{The zero-knowledge direction}
\label{sec:crypto-zk}

The natural next question for a privacy-sensitive counterparty is: can I prove
``my catalogue contains something under $\mathrm{weight}(..5\mathrm{kg})$ and
$\mathrm{location}(\mathrm{EU})$'' \emph{without revealing the catalogue}?

This is the zero-knowledge question in its original sense
\citep{goldwasser1989}: convince a verifier that a statement holds while
revealing nothing beyond its truth. It is heavy research --- succinct proofs
over trie walks --- and \odag{} deliberately walls it until a real consumer appears. But the positioning is
worth recording, because \odag{} is unusually well suited when the time comes:

\begin{itemize}[leftmargin=1.4em]
\item deterministic canonical encoding: the statement being proved is about a
  well-defined object, not about one serialisation among many;
\item one query primitive: a single circuit shape covers the whole query
  language;
\item \textbf{no floating point anywhere} --- values are pairs of integers.
  Arithmetic circuits are built over finite fields, where rationals are
  natural and floats are a nightmare. \odag{}'s exactness commitment
  (\S\ref{sec:keep}), adopted for canonicality, turns out to be exactly the
  property a proof system wants.
\end{itemize}

That last point is a good illustration of a theme in this paper: the refusals
compound. Refusing floats to protect canonical roots also made the structure
circuit-friendly, years before anyone asked.

\subsection{Governance: canonical forms, plural}

A final point, less technical but arguably the most important for a
decentralised deployment.

Every attempt at a single official classification of the world has failed, and
Borges gave the argument its definitive form: since we do not know what the
universe is, every classification of it is arbitrary and conjectural
\citep{borges1942}. Foucault made the historical version of the same point ---
classification schemes are not refined over time so much as replaced.

The constructive reading, for a system design, is that \emph{no single
taxonomy deserves to be baked into identity}. \odag{}'s answer is
architectural rather than editorial:

\begin{itemize}[leftmargin=1.4em]
\item there is no official store --- anyone publishes a root;
\item vocabularies travel \emph{inside} stores as ordinary declarations, so
  adopting a vocabulary is merging a published root, not installing software;
\item adoption is by explicit reference to a root, so what you adopted is
  exactly what you can name and check;
\item disagreement is representable: two communities may hold divergent
  overlays of a shared core and both be well-formed.
\end{itemize}

Wilkins wanted the one true tree. The answer here is closer to Borges:
canonical \emph{forms}, plural --- one per community, per purpose, per moment
--- each canonical in itself, each mergeable with the others, and none final.

\subsection{Honest limits, again}

\begin{itemize}[leftmargin=1.4em]
\item \textbf{L1: proofs attest structure, never truth.} No amount of
  cryptography makes an assertion true. Attribution and bonds change the
  incentives, not the epistemology.
\item \textbf{L2: everything is relative to a pinned interpretation context.}
  A proof about a typed value must pin the declarations and the registry
  version. A verifier running different arithmetic must refuse rather than
  misinterpret --- and \odag{}'s certificates do refuse.
\item \textbf{Grow-only merge} remains the standing constraint on any
  multi-writer deployment.
\item \textbf{The infrastructure wall is real.} Running a node, funding it,
  and buying storage rent is a genuine barrier --- measured in this project's
  own logs at minutes-to-hours for a first-time user, not seconds. The
  practical response has been to make the interesting properties available
  \emph{before} the network step: a local content-addressed store gives
  canonical roots, snapshots, certificates and merge with no node at all, so
  the ideas can be evaluated before the infrastructure is confronted.
\end{itemize}
